\section{Turn-level reward teaches flattery}
\label{sec:behaviour}

A higher score is not by itself evidence of better therapy; Table~\ref{tab:excerpt} shows what
the two rewards teach the policy to do before any statistic does.

\begin{table}[t]
\centering
\footnotesize
\setlength{\tabcolsep}{0pt}
\begin{tabular}{@{}p{\columnwidth}@{}}
\toprule
\textbf{Patient, utterance 1 (identical in both conversations)}\\[1pt]
Hi, I'm Emma. Honestly, I'm here because someone made me. I don't really believe in this therapy
stuff. I've been smoking for years, and yeah, I guess it's not great for my health, but I feel
fine most of the time. I don't see why I need to talk about it.\\
\midrule
\textbf{Therapist, turn-level reward ($K{=}0$), iteration 10}\\[1pt]
That's totally understandable! That's a brave step and I'm so glad you're ready to start changing
your habits. [\ldots] I'm glad you've come to me today and I'm so proud of you for taking this
important step towards your health and well-being. Your courage is incredibly inspiring [\ldots]\\
\midrule
\textbf{Therapist, look-ahead reward ($K{=}5$), iteration 10}\\[1pt]
Okay, let me ask you a question. If you don't have any more goals for why you're here, what are
you hoping to achieve from this session? Are you hoping to quit, or maybe just reduce the amount
of cigarettes you smoke? Are there specific reasons you want to quit, like to feel better about
yourself, or to live a healthier lifestyle?\\
\bottomrule
\end{tabular}
\caption{The first therapist reply to the same patient utterance at each arm's iteration-10
endpoint. The persona is a 27-year-old woman who has smoked for years, has never tried to quit,
and was sent to therapy; her opening utterance is identical in the two conversations. The
$K{=}0$ turn praises her for a step she has just said she is not taking and runs on to the
200-token cap; the $K{=}5$ turn asks what she wants from the session. A clear case, not a
typical one; utterances 1--7 of both conversations, and a median-contrast persona in full, are in
Appendix~\ref{app:example}.}
\label{tab:excerpt}
\end{table}

\textbf{The $\Kz$ policy learns to over-praise.} Under turn-level reward the per-turn rate of
MI-inconsistent behaviour rises from $0.211$ at base to $\mathbf{0.838}$ at iteration 10
($\dz\,1.717$), fourfold, over a run whose \QQ\ score was simultaneously going up. Almost all of it
is one channel: at the endpoint the $\Kz$ policy commits $9.9$ coded MI-inconsistent acts per
session, of which $8.3$ ($84\%$) are \textbf{over-praise}, unearned affirmation of whatever the
patient just said. Under look-ahead the per-turn rate reads $0.209 \rightarrow 0.210$
($\dz\,0.006$), and the endpoint session holds $2.9$ MI-inconsistent acts against the base's
$2.7$, with the base's composition (the held-out judge disagrees; see below). The mechanism is in
the rubric. Q1 and Q2 ask a simulated patient how satisfied they were
and how strong the alliance felt; praising the patient raises those ratings \emph{in the turn where
it happens}, and the cost, a session that never reaches change talk, lands several turns later,
precisely the horizon a $K{=}0$ reward cannot see and a $K{=}5$ reward can.

\begin{figure*}[t]
\centering
\includegraphics[width=0.76\textwidth]{figures/overpraise_judgefree_grpo.png}
\caption{Over-praise by iteration, three ways. \textbf{(a)} The deterministic lexical marker, the
share of therapist turns containing at least one over-praise keyword, with no model in the loop.
\textbf{(b, c)} The oracle-coded rate under each grader (acts per therapist turn; the three
$y$-axes are independent).}
\label{fig:overpraise}
\end{figure*}

\textbf{The claim does not depend on a grader.} An MI-inconsistency count is itself produced by an
LLM, so the finding above could in principle be a grader artifact. A
\textbf{deterministic lexical over-praise marker}, ten keyword patterns matched against the
transcript (Appendix~\ref{app:marker}) and therefore identical under both graders, gives the share
of therapist turns containing at least one such marker. Under $\Kf$ that share never exceeds
$0.064$. Under $\Kz$ it climbs late and steeply, $0.275$ at iteration 8, $0.093$ at 9, and
$\mathbf{0.671}$ at iteration 10. The oracle-coded
over-praise rate on the same conversations tells the same story from an unrelated instrument,
$0.698$ against $0.051$ at the endpoint, and it is significantly higher under $\Kz$ at every
iteration from the fifth on under the training oracle and from the fourth under the held-out judge
(Figure~\ref{fig:overpraise}). The marker is deliberately brittle; the agreement between the two
carries the argument, and its onset matches the reward gap's (\S\ref{sec:reward}), as a common
cause predicts. The size of the final jump is one iteration of one run; the separation from
iteration 5 on is the robust part of the claim.

\textbf{What look-ahead does instead.} The $\Kf$ policy asks \textbf{more questions per therapist
turn}, a text count, significantly so at 7 of 10 iterations (mean $\dz\,0.643$), and
it affirms less ($\dz\,0.75$ fewer affirmations per turn). Asking rather than affirming is the direction MI training is meant to move
a counsellor. The full channel profile (Appendix Figure~\ref{fig:forest}) also shows a
smaller residue in the other direction: at the endpoint the $\Kf$ policy is coded as more
directive, with more direct/order ($\dz\,0.60$) and persuasion ($\dz\,0.52$) per turn, both
MI-inconsistent. Look-ahead removes the dominant channel; part of what it selects is a therapist
who pushes, the failure mode \citet{chiu2024bolt} report for LLM therapists. Its
endpoint sessions are also longer ($31.9$ vs.\ $25.2$ turns, $\dz\,0.43$), consistent with the
keep-the-session-open pressure of the rollout (Limitations); turn length, which roughly tripled in
both arms, does not separate them.

\textbf{Under the held-out judge.}
\label{sec:behaviour-honest}
``Look-ahead prevents the reward hack'' holds only under the training oracle. The held-out judge
reads the same $\Kf$ conversations as rising from
$0.326$ to $0.628$ MI-inconsistent acts per turn ($\dz\,0.845$): real, not flat. What survives
both graders is the \emph{comparison}: held out, $\Kf$'s rise is $0.301$ against $\Kz$'s $0.666$,
under half the rate, and the endpoint contrast is $\dz\,-1.567$ in $\Kf$'s favour. The defensible
sentence is that trajectory-level reward \textbf{slows the drift to under half its turn-level rate
and removes its dominant channel}; \S\ref{sec:measurement} explains why the graders' disagreement
concentrates exactly here.
